\section{Final Report}

The final report is the main written record of your project. It should present the full product development process in a self-contained way so that a reader can understand the problem, the evidence gathered, the design decisions made, and the proposed final concept without needing to read your weekly submissions first.

Submit one group report in PDF format through ITEL. The cover page should identify the group members, project title, and any other details requested by the lecturer. The final report should normally be about 20 to 25 pages long

\subsection{Recommended Format and Tools}

You may use \LaTeX{} or another suitable word processor (e.g., Microsoft Word). However, \LaTeX{} is recommended because it supports consistent formatting, references, tables, and figure management, but it is not compulsory if your group can produce an equally professional report in another format.

Regardless of the software used, the report should be:
\begin{itemize}
    \item clearly structured with section headings;
    \item written in formal academic English;
    \item supported by captions, references, and labelled figures or tables; and
    \item checked carefully for consistency, grammar, and technical accuracy.
\end{itemize}

\subsection{Suggested Report Structure}

The structure below is recommended because it aligns with the weekly APID workflow and the common grading criteria.


\begin{enumerate}
    \item \textbf{Report Title}
    \begin{enumerate}
        \item Provide a clear, concise title accurately reflecting the content of the report.
    \end{enumerate}

    \item \textbf{Introduction}
    \begin{enumerate}
        \item The introduction should concisely and comprehensively outline the issue your product aims to address and explain why the product is necessary.
        \item Provide an overview of the main sections of the report to clearly indicate the content that follows.
    \end{enumerate}

    \item \textbf{Opportunity Identification}
    \begin{enumerate}
        \item Develop your group’s \textbf{\textcolor{red}{Innovation Charter}} as a single clear and concise statement. The charter should follow the required guideline, provide a strong sense of direction for product development, and be appropriately balanced. A strong charter is specific enough to guide decision-making, but not so narrow that it limits creativity or so broad that it becomes vague and unfocused.
        \item Provide a comprehensive and insightful \textbf{\textcolor{red}{explanation}} of why an innovation charter is important in the product development process. Explain how it helps guide product decisions, focus innovation efforts, and encourage creativity within an organization. A strong response should include multiple clear and relevant examples showing how an innovation charter can shape product development and improve innovation outcomes.
        \item List all product or innovation \textbf{\textcolor{red}{opportunities}} identified from your group’s innovation charter. Each opportunity should be clearly connected to the charter and demonstrate strong relevance to the intended innovation direction. Strong opportunities should reflect a clear understanding of customer desirability, business viability, and technological feasibility.
        \item Propose at least \textbf{\textcolor{red}{five high-quality opportunities}}. The opportunities should be clearly written, relevant, realistic, and meaningful in potential impact. They should also be well aligned with the innovation charter and show a strong balance across user value, business potential, and technical feasibility.
    \end{enumerate}

    \item \textbf{Product Planning}
    \begin{enumerate}
        \item Since the opportunity funnel may initially capture numerous ideas, many may not align well with your firm's goals or capabilities. It is therefore crucial to identify and select the most promising opportunities. To manage your project portfolio effectively, evaluate competitive strategies, market segmentation, and technological trajectories. Specifically, address the following:
        \begin{enumerate}
            \item Identify and evaluate the relevant \textbf{\textcolor{red}{market segment(s)}} for the proposed product. Explain the needs, preferences, and characteristics of each segment, and discuss how these influence product design and feature selection.
            \item Identify any new or \textbf{\textcolor{red}{emerging technologies}} that should be incorporated into the product. Explain how these technologies enhance product functionality, performance, or customer value, and justify their adoption in relation to technological trends or trajectories.
            \item Define the key \textbf{\textcolor{red}{manufacturing and service objectives}} for the product, as well as the major constraints that may affect production, delivery, or support. Discuss how these constraints can be managed or addressed.
            \item Clearly specify the primary \textbf{\textcolor{red}{target market(s)}} and propose a target selling price. Justify the price based on the characteristics, needs, and purchasing power of the selected market segment(s).
            \item State the main \textbf{\textcolor{red}{assumptions and constraints}} underlying the planning process. Identify the key stakeholders involved in product development and briefly explain their roles and contributions.
        \end{enumerate}


        \item Based on these analyses, create a \textbf{\textcolor{red}{mission statement}} that includes the following details:
        \begin{enumerate}
            \item Product description
            \item Benefit proposition
            \item Key business goals
            \item Primary market
            \item Secondary markets
            \item Assumptions and constraints
            \item Stakeholders
        \end{enumerate}
    \end{enumerate}

    \item \textbf{Development Processes and Organization}
    \begin{enumerate}
        \item Propose the most appropriate \textbf{\textcolor{red}{product development process}} for your product, such as the Generic Product Development Process, Spiral Product Development Process, or Complex System Development Process. Briefly justify your selection by explaining why the chosen process is suitable for the nature of the product, the level of uncertainty involved, the complexity of development, and the needs of the project. A strong response should demonstrate a clear understanding of the selected process and how it supports effective product development.
        \item Define the optimal \textbf{\textcolor{red}{process flow}} for developing the product. Present the key stages or activities in a logical sequence and show how the development effort should progress from idea generation to product launch or implementation. The process flow should be clear, well organized, and appropriate for the selected development process.
        \item Suggest the most suitable \textbf{\textcolor{red}{organizational structure}} to support the product development effort. Briefly justify your choice by explaining how the structure supports coordination, communication, workflow efficiency, and alignment with organizational goals and business strategy. A strong response should clearly show why the proposed structure is appropriate for the product and development context.
    \end{enumerate}
    
    % \item \textbf{Optional Elements}
    
    % \begin{enumerate}
    %     \item Concept Generation
    %     \item Concept Selection
    %     \item Concept Testing
    %     \item Intellectual Property Management
   
    %     Although optional, it is highly recommended that your group includes summaries of consensus from previous activities such as. Even though these elements have been previously assessed through multimedia presentations or practical activities, documenting them here ensures comprehensive coverage of your group's findings. You may directly copy related materials (e.g., concept sketches from Microsoft Whiteboard) without extensive rewriting.
    % \end{enumerate}

\end{enumerate}

\subsection{Optional Elements}
These sections are optional because they have been previously assessed through multimedia presentations or practical activities. However, including summaries of consensus from these activities can ensure comprehensive coverage of your group's findings. You may directly copy related materials (e.g., concept sketches from Microsoft Whiteboard) without extensive rewriting.




Table~\ref{tab:session0-final-report-needs} shows a simplified example of how customer needs can be documented.

\begin{table}[htbp]
\centering
\caption{Example format for translating customer statements into interpreted needs}
\label{tab:session0-final-report-needs}
\renewcommand{\arraystretch}{1.2}
\begin{tabularx}{\textwidth}{|>{\raggedright\arraybackslash}p{0.24\textwidth}|>{\raggedright\arraybackslash}X|>{\raggedright\arraybackslash}X|}
\hline
\textbf{Prompt or theme} & \textbf{Customer statement} & \textbf{Interpreted need} \\
\hline
Typical use & ``I need to check the nutrient tank too often.'' & The system should make nutrient-level monitoring quick and easy. \\
\hline
Dislike & ``Water flow becomes inconsistent when the pump clogs.'' & The system should maintain stable nutrient flow and allow easy maintenance. \\
\hline
Suggested improvement & ``I want to know immediately when the pH changes.'' & The system should provide timely pH monitoring and alert functions. \\
\hline
\end{tabularx}
\end{table}

Table~\ref{tab:session0-final-report-metrics} gives a matching example for the specification stage.

\begin{table}[htbp]
\centering
\caption{Example format for linking needs to measurable metrics}
\label{tab:session0-final-report-metrics}
\renewcommand{\arraystretch}{1.2}
\begin{tabularx}{\textwidth}{|>{\raggedright\arraybackslash}X|>{\raggedright\arraybackslash}p{0.2\textwidth}|>{\raggedright\arraybackslash}p{0.12\textwidth}|>{\raggedright\arraybackslash}X|}
\hline
\textbf{Customer need} & \textbf{Metric} & \textbf{Unit} & \textbf{Justification} \\
\hline
Stable nutrient delivery & Nutrient flow-rate variation & L/min & Indicates whether the delivery system remains consistent during operation. \\
\hline
Easy monitoring & Display response time after sensor update & s & Shows how quickly a user receives useful system feedback. \\
\hline
Safe operation & Leakage rate during 24-hour operation & mL/h & Relates directly to reliability, cleanliness, and maintenance burden. \\
\hline
\end{tabularx}
\end{table}

\subsection{Submission Requirements}

Before submission, confirm that the report:
\begin{itemize}
    \item is complete and internally consistent;
    \item uses captions and references correctly;
    \item cites external sources where needed;
    \item follows any naming convention announced on ITEL; and
    \item reflects your own checked and edited understanding, including when LLM tools were used.
\end{itemize}

\subsection{About the Use of LLMs in the Report}

LLMs may help you improve structure, grammar, or brainstorming efficiency, but they cannot replace technical judgement. If you use an LLM while preparing the final report, verify every claim, remove unsupported content, and rewrite the material so that it matches your project evidence and the actual task requirements. Weak reports often result from unedited LLM output that sounds polished but does not answer the question accurately.
